% =====================================================================
% Resumo, palavras-chave e sumário. Sem número de seção.
% Resumo: parágrafo único, 150 a 300 palavras (roteiro, seção 4).
% =====================================================================
\section*{Resumo}
\addcontentsline{toc}{section}{Resumo}

Este trabalho relaciona a história do cálculo numérico com a ciência da computação e aprofunda os métodos de passo único para equações diferenciais ordinárias. A parte histórica vai das tabelas e dos instrumentos de cálculo até os computadores eletrônicos, a aritmética de ponto flutuante e as linguagens científicas. O problema de engenharia é o circuito RLC série, usado como modelo de um filtro ressonante de telecomunicações. O circuito é descrito por duas equações de estado, para a corrente e a tensão no capacitor. Essas equações são resolvidas pelos métodos de Euler, de Euler melhorado e de Runge-Kutta de quarta ordem, escritos à mão em Scilab. As soluções numéricas são comparadas com as soluções exatas do degrau, nos regimes subamortecido, crítico e sobreamortecido, e da senoide, dentro e fora da ressonância. Também são comparadas com o comando ode do Scilab. O estudo varia o passo e o fator de qualidade do circuito. As ordens observadas ficam em 1, 2 e 4, como prevê a teoria. O limite de estabilidade do método de Euler é igual ao produto da resistência pela capacitância: 20 microssegundos com fator de qualidade 5 e 2 microssegundos com fator de qualidade 50. Nesse segundo caso, o Runge-Kutta de quarta ordem continua estável até cerca de 285 microssegundos. Para um erro de um microampère, ele precisa de cerca de 1\,300 avaliações da função, contra 3,3 milhões do método de Euler. Conclui-se que a ordem reduz o custo de uma precisão desejada e que a região de estabilidade decide se a simulação de um filtro seletivo continua fiel ao circuito.

\vspace{0.6em}
\noindent\textbf{Palavras-chave:} cálculo numérico; método de Runge-Kutta; circuito RLC; estabilidade numérica; Scilab.

\clearpage
\tableofcontents
\clearpage
